% ============================================================
% 纯答案册·课时06-两条直线的位置关系 —— 工具/答案抽册器.py 抽册生成件（勿手改；第三档＝纯答案单独成册）
% 源件（只读）：C:/提示词/工作区/M3-第2章量产0913/成卷/导学件/课时06-两条直线的位置关系/main.tex（md5 7e88aff887d7486a7959c9bfc4fb4081）
%% 口径：题面不重印；逐题＝题号(印面号同正册)＋[答案]＋[详解]，值/详解照源件逐字
%       （钉值门硬断言：册值≡正册件内值≡值台账值tex，逐字全等）。册序＝正册装配序＝印面号 1..21 升序（同序同号）；键序断言基准＝题面库冻结 manifest canonical 键序。
%% 三档导出：整体 main-true／纯题 main-false（在产）｜本册＝纯答案册（本件）
%% 双档：ansbook-true.tex＝详解印本；ansbook-false.tex＝纯值速查本（详解不印）
% 编译：xelatex <壳> ×2，cwd＝本目录；sty＝qp-m3.sty 本地挂载（md5 7c3930362be8a0a2bdf21bbf8ac16573，锁校验）
% ============================================================
\documentclass[fontset=none]{ctexart}
\usepackage{qp-m3}
% —— 印面逐字符号路由补钉（承源件；− 走 TNR、▱ NSC 子块；禁改 sty 保 md5） ——
\xeCJKDeclareCharClass{Default}{"2212}%
\setCJKfamilyfont{nscblk}[Path=C:/提示词/工作区/字替对照-0909/variantF/fonts/,BoldFont=NSC-w600.ttf]{NSC-w500.ttf}%
\xeCJKDeclareSubCJKBlock{nscblk}{"25B1}%
\newcommand{\pxparallelogram}{{\CJKfamily{nscblk}▱}}%
% —— 断行微调（承母版实测档，三零门承重；\emergencystretch 不另设） ——
\xeCJKsetup{CJKglue={\hskip 0.10em plus 0.08em minus 0.05em}}%
% —— 纯答案册全线括线（长详解可跨栏断，承练习件全线括线口径；灰底盒不可跨栏断） ——
\ansblockgrayfalse
% —— 双档开关（false 档＝纯值速查：答案印、详解不印；件面开关，不动 sty 本体） ——
\newif\ifansbookdetail
\ifdefined\ansbookdetail
  \ifnum\ansbookdetail=0 \ansbookdetailfalse\else\ansbookdetailtrue\fi
\else
  \ansbookdetailtrue
\fi
% —— 页脚件名（承源件章名；件名＝<件型>答案册） ——
\renewcommand{\qpzhangming}{第二章\quad 平面解析几何}%
\renewcommand{\qpjianming}{导学件答案册}%

\begin{document}
\zhangtitle{第二章\quad 平面解析几何}
\jietitle{2.2.3 两条直线的位置关系}
\par\glueguard{1}\addvspace{2pt}\noindent\biaoqian{【纯答案册】}{\fangsong 题面不重印\quad 印面号与正册同号同序}\par\addvspace{2pt}

\begin{multicols}{2}
\emergencystretch=1em

\begin{ansblock}[2章-练-课时06-T2]
% ans:2章-练-课时06-T2
\ansitem{1}{B}
\ifansbookdetail
\ansnote{详解}{①\(k_{AB}=(8-2)/(4-1)=2=k_1\)，但只知斜率相等，\(l_1\)与\(l_2\)可能重合，①不入选；\par\noindent ②\(k_{PQ}=0\)，\(l_1\)为水平线；\(l_2\)平行于\(x\)轴且不过\(P\)，两水平线不重合，平行，②入选；\par\noindent ③\(k_1=(-2-0)/(-5+1)=1/2\)，\(k_2=(5-3)/(0+4)=1/2\)；检验\(M(-1,0)\)不在\(l_2\)（\(l_2: y=x/2+5\)，\(x=-1\)时\(y=9/2\)）上，故平行，③入选．选：B．}
\fi
\end{ansblock}

\begin{ansblock}[2章-练-课时06-T3]
% ans:2章-练-课时06-T3
\ansitem{2}{D}
\ifansbookdetail
\ansnote{详解}{\(k_1=(-3-0)/(1-0)=-3\)；\(l_2: ax+y-2=0\)的斜率\(k_2=-a\)．\(l_1\parallel l_2\iff k_1=k_2\)，得\(-a=-3\)，\(a=3\)．重合检验：\(a=3\)时\(l_2: 3x+y-2=0\)，取\(l_1\)上点\((0,0)\)代入得\(-2\neq0\)，\((0,0)\)不在\(l_2\)上，两直线不重合，平行成立．选：D．}
\fi
\end{ansblock}

\begin{ansblock}[2章-练-课时06-T4]
% ans:2章-练-课时06-T4
\ansitem{3}{C}
\ifansbookdetail
\ansnote{详解}{\(k_1=\tan60^\circ=\sqrt{3}\)；\(k_2=(2\sqrt{3}-\sqrt{3})/(-2-1)=-\sqrt{3}/3\)．\(k_1\cdot k_2=\sqrt{3}\times(-\sqrt{3}/3)=-1\)，且\(k_1\neq k_2\)（两线不平行、更不重合），故\(l_1\perp l_2\)．选：C．}
\fi
\end{ansblock}

\begin{ansblock}[2章-练-课时06-T6]
% ans:2章-练-课时06-T6
\ansitem{4}{C}
\ifansbookdetail
\ansnote{详解}{垂直\(\iff k(k-1)+(1-k)(2k+3)=0\)．展开：\(k^2-k+2k+3-2k^2-3k=-k^2-2k+3=0\)，即\(k^2+2k-3=0\)，解得\(k=1\)或\(k=-3\)．选：C．}
\fi
\end{ansblock}

\begin{ansblock}[2章-练-课时06-T11]
% ans:2章-练-课时06-T11
\ansitem{5}{−1}
\ifansbookdetail
\ansnote{详解}{\(l_1\parallel l_2\iff 2\times1-2m=0\)，得\(m=1\)（此时\(l_1: 2x+y+2=0\)与\(l_2\)不重合）；\(l_1\perp l_3\iff 2\times1+n=0\)，得\(n=-2\)．故\(m+n=-1\)．}
\fi
\end{ansblock}

\begin{ansblock}[2章-练-课时06-T14]
% ans:2章-练-课时06-T14
\ansitem{6}{(1) m=1 或 6；(2) m=3 或 −4}
\ifansbookdetail
\ansnote{详解}{\(k_1=(2-m)/(m-4)\)，\(k_2=(m+2-2)/(-2-1)=-m/3\)．\par\noindent (1) \(k_1=k_2\)：\((2-m)/(m-4)=-m/3\)，即\(3(2-m)=-m(m-4)\)，\(m^2-7m+6=0\)，\(m=1\)或\(6\)．检验不重合：\(m=1\)时\(l_1\)过\((3,1)\)、\((0,2)\)，斜率为\(-1/3\)，\(C(1,2)\)不在\(l_1\)上；\(m=6\)时\(l_1: y=-2x+12\)，\(C(1,2)\)不在其上，均不重合，故\(m=1\)或\(6\)；\par\noindent (2) \(m=0\)时\(k_2=0\)（水平），需\(l_1\)铅直，但\(A(3,0)\)、\(B(-1,2)\)横坐标不同，斜率存在，舍；\(k_2\neq0\)时\(k_1\cdot k_2=-1\)：\((2-m)/(m-4)\cdot(-m/3)=-1\)，即\(m(2-m)=-3(m-4)\)，\(m^2-m-12=0\)，\(m=3\)或\(-4\)．验：\(m=3\)时\(k_1=1\)，\(k_2=-1\)；\(m=-4\)时\(k_1=-3/4\)，\(k_2=4/3\)，乘积均为\(-1\)，故\(m=3\)或\(-4\)．}
\fi
\end{ansblock}

\begin{ansblock}[2章-练-课时06-T1]
% ans:2章-练-课时06-T1
\ansitem{7}{B}
\ifansbookdetail
\ansnote{详解}{\(A_1A_2+B_1B_2=2a\times1+1\times[-(a+1)]=a-1=0\)，得\(a=1\)（\(a=-1\)时\(l_2: x+2=0\)铅直而\(l_1\)斜率为2，不垂直，已排除）．此时\(l_1: 2x+y-2=0\)，\(l_2: x-2y+2=0\)．由\(y=2-2x\)代入：\(x-2(2-2x)+2=0\)，\(5x=2\)，\(x=2/5\)，\(y=6/5\)．交点\((2/5, 6/5)\)，选：B．}
\fi
\end{ansblock}

\begin{ansblock}[2章-练-课时06-T5]
% ans:2章-练-课时06-T5
\ansitem{8}{A}
\ifansbookdetail
\ansnote{详解}{\(k_{BC}=2/(-8)=-1/4\)，则\(k_{AH}=4\)；\(k_{BH}=1/(-5)=-1/5\)，则\(k_{AC}=5\)．设\(A(x,y)\)：\((y-2)/(x+3)=4\)且\((y-3)/(x+6)=5\)，即\(y=4x+14\)且\(y=5x+33\)，解得\(x=-19\)，\(y=-62\)．验第三组垂直：\(k_{CH}=(2-3)/(-3+6)=-1/3\)，而\(k_{AB}=(1+62)/(2+19)=3\)，乘积为\(-1\)，成立．选：A．}
\fi
\end{ansblock}

\begin{ansblock}[2章-练-课时06-T7]
% ans:2章-练-课时06-T7
\ansitem{9}{B}
\ifansbookdetail
\ansnote{详解}{垂直\(\iff a\cdot a+(b+2)(b-2)=0\)，即\(a^2+b^2=4\)．由\(a^2+b^2\geqslant2|ab|\)得\(|ab|\leqslant2\)；取\(a=b=\sqrt{2}\)（满足\(a^2+b^2=4\)，此时两直线为\(\sqrt{2}x+(2+\sqrt{2})y+4=0\)与\(\sqrt{2}x+(\sqrt{2}-2)y-3=0\)，\(A_1A_2+B_1B_2=2+(2+\sqrt{2})(\sqrt{2}-2)=2+(2-4)=0\)，垂直成立），\(ab=2\)．故\(ab\)的最大值为2，选：B．}
\fi
\end{ansblock}

\begin{ansblock}[2章-练-课时06-T9]
% ans:2章-练-课时06-T9
\ansitem{10}{B}
\ifansbookdetail
\ansnote{详解}{联立\(2x+y-3=0\)与\(x-3y+2=0\)：由第一式\(y=3-2x\)代入第二式：\(x-3(3-2x)+2=0\)，得\(x=1\)，\(y=1\)，交点\((1,1)\)．\(k_1=-2\)，所求直线斜率\(k=1/2\)．\(y-1=(1/2)(x-1)\)，即\(x-2y+1=0\)，选：B．}
\fi
\end{ansblock}

\begin{ansblock}[2章-练-课时06-T12]
% ans:2章-练-课时06-T12
\ansitem{11}{1}
\ifansbookdetail
\ansnote{详解}{四边形由\(l_1\)、\(l_2\)与\(x\)轴、\(y\)轴围成，坐标轴交点处的内角为\(90^\circ\)，四边形有外接圆\(\iff\)对角互补\(\iff l_1\)与\(l_2\)的交角为\(90^\circ\)，即\(l_1\perp l_2\)．\(k_1=-1/3\)，\(k_2=3k\)，垂直\(\iff(-1/3)\times3k=-1\)，得\(k=1\)．（检验：\(l_1\)不过原点、不平行于坐标轴；\(k=1\)时\(l_2: 3x-y+1=0\)同样，四边形存在且对角互补，有外接圆．）}
\fi
\end{ansblock}

\begin{ansblock}[2章-练-课时06-T15]
% ans:2章-练-课时06-T15
\ansitem{12}{(1) D(−1,6)；(2) 是菱形}
\ifansbookdetail
\ansnote{详解}{(1) 平行四边形对角线互相平分：\(AC\)中点与\(BD\)中点重合．设\(D(x,y)\)：\(((1+3)/2, (2+4)/2)=((5+x)/2, (0+y)/2)\)，得\(x=-1\)，\(y=6\)，即\(D(-1,6)\)；\par\noindent (2) 邻边相等即菱形：\(|AB|^2=16+4=20\)，\(|BC|^2=4+16=20\)，\(|AB|=|BC|\)，故\pxparallelogram{}ABCD 为菱形．（或对角线：\(k_{AC}=1\)，\(k_{BD}=6/(-6)=-1\)，\(AC\perp BD\)且互相平分，也为菱形．）}
\fi
\end{ansblock}

\begin{ansblock}[2章-练-课时06-T13]
% ans:2章-练-课时06-T13
\ansitem{13}{(1,1)；√10}
\ifansbookdetail
\ansnote{详解}{按\(\lambda\)整理：\((x+y-2)+\lambda(3x+2y-5)=0\)，令两因式同时为零：\(x+y-2=0\)且\(3x+2y-5=0\)，解得\(x=y=1\)，定点\(Q(1,1)\)．\(\lambda\)变化时\(l\)绕\(Q\)转动（斜率可取除\(-3/2\)外的一切值，且\(\lambda=-1/2\)时为铅直线\(x=1\)，方向覆盖齐全）．\(d\)最大\(\iff l\perp PQ\)，最大距离\(|PQ|=\sqrt{(-2-1)^2+(0-1)^2}=\sqrt{10}\)．}
\fi
\end{ansblock}

\begin{ansblock}[2章-练-课时06-T8]
% ans:2章-练-课时06-T8
\ansitem{14}{ACD}
\ifansbookdetail
\ansnote{详解}{A：\(l_2\)按\(a\)整理：\(a(x-2y)+(3y-1)=0\)，令\(x-2y=0\)且\(3y-1=0\)，得\((2/3, 1/3)\)，恒过，A对；\par\noindent B：\(a=1\)时\(l_1: x+y-1=0\)，\(l_2: x+y-1=0\)，重合非平行；\(a=-3\)时\(k_1=-1/3\)，\(l_2: -3x+9y-1=0\)的斜率为\(1/3\)，不平行，B错；\par\noindent C：\(A_1A_2+B_1B_2=a+a(3-2a)=2a(2-a)=0\)，得\(a=0\)或\(2\)，C对；\par\noindent D：\(a>0\)时\(l_1: y=-x/a+1\)，斜率为负、纵截距\(1>0\)、横截距\(a>0\)，图象过一、二、四象限，不过第三象限，D对．选：ACD．}
\fi
\end{ansblock}

\begin{ansblock}[2章-练-课时06-T10]
% ans:2章-练-课时06-T10
\ansitem{15}{BCD}
\ifansbookdetail
\ansnote{详解}{\(AB\)中点\((-2,2)\)，\(k_{AB}=1\)，中垂线\(y-2=-(x+2)\)，即\(x+y=0\)．与欧拉线联立：\(x=-1\)，\(y=1\)，外心\(O'(-1,1)\)，\(R^2=9+1=10\)．\par\noindent 设\(C(m,n)\)：重心\(((m-4)/3, (n+4)/3)\)在欧拉线上，得\(m-n-2=0\)；\(|O'C|^2=10\)，即\((m+1)^2+(n-1)^2=10\)，展开\(m^2+n^2+2m-2n=8\)．代入\(n=m-2\)：\(2m^2-4m=0\)，\(m=0\)或\(2\)，\(C(2,0)\)或\((0,-2)\)，C对；\par\noindent 重心为\((-2/3, 4/3)\)或\((-4/3, 2/3)\)，A所给数值错，A错；\par\noindent 垂心：由欧拉关系\(H=3G-2O'\)（\(G\)为重心）：\(C(2,0)\)时\(H=3(-2/3, 4/3)-2(-1,1)=(0,2)\)；\(C(0,-2)\)时\(H=3(-4/3, 2/3)-2(-1,1)=(-2,0)\)．（直接验证\(C(2,0)\)情形：\(AC\)水平，过\(B\)的高为\(x=0\)；过\(C\)垂直\(AB\)的高为\(y=-x+2\)，交点\((0,2)\)．）B对；\par\noindent 面积比：\(AB\)即\(x-y+4=0\)与欧拉线平行，\(C(2,0)\)到欧拉线距离\(4/\sqrt{2}\)，到\(AB\)距离\(6/\sqrt{2}\)，相似比\(2/3\)，小三角形与大三角形面积比\(4/9\)，故两部分之比\(4:5\)，即\(4/5\)，D对．选：BCD．}
\fi
\end{ansblock}

\begin{ansblock}[2章-练-课时06-T16]
% ans:2章-练-课时06-T16
\ansitem{16}{(1) 5x+y−13=0；(2) 7/2}
\ifansbookdetail
\ansnote{详解}{(1) \(k_{BC}=(-2-3)/(3-2)=-5\)，\(y-3=-5(x-2)\)，即\(5x+y-13=0\)；\par\noindent (2) \(|BC|=\sqrt{1+25}=\sqrt{26}\)，\(A\)到\(BC\)的距离\(d=|5\times1+1-13|/\sqrt{26}=7/\sqrt{26}\)，故\(S=(1/2)\times\sqrt{26}\times(7/\sqrt{26})=7/2\)．}
\fi
\end{ansblock}

\begin{ansblock}[2章-导-课时06-G1]
% ans:2章-导-课时06-G1
\ansitem{17}{D}
\ifansbookdetail
\ansnote{详解}{\(k_1=\tan135^\circ=-1\)；\(k_2=(-6-(-1))/(3-(-2))=-5/5=-1\)．两斜率相等，但仅凭斜率相等不能排除重合（题面未给两点异于\(l_1\)的信息），故\(l_1\)与\(l_2\)平行或重合．选：D．}
\fi
\end{ansblock}

\begin{ansblock}[2章-导-课时06-G2]
% ans:2章-导-课时06-G2
\ansitem{18}{B}
\ifansbookdetail
\ansnote{详解}{平行\(\iff a(a-1)-2\times1=0\)，即\(a^2-a-2=0\)，\(a=2\)或\(a=-1\)．检验：\(a=-1\)时两直线分别为\(-x+2y-1=0\)与\(x-2y+1=0\)，系数成比例，重合，舍；\(a=2\)时两直线\(2x+2y-1=0\)与\(x+y+4=0\)平行不重合．故\(a=2\)，选：B．}
\fi
\end{ansblock}

\begin{ansblock}[2章-导-课时06-G3]
% ans:2章-导-课时06-G3
\ansitem{19}{C}
\ifansbookdetail
\ansnote{详解}{联立\(y=2x+10\)与\(y=x+1\)：\(x=-9\)，\(y=-8\)，交点\((-9,-8)\)．代入\(y=ax-2\)：\(-8=-9a-2\)，得\(a=2/3\)，选：C．}
\fi
\end{ansblock}

\begin{ansblock}[2章-导-课时06-G5]
% ans:2章-导-课时06-G5
\ansitem{20}{3x+y+1=0}
\ifansbookdetail
\ansnote{详解}{设\(l\)与\(l_1\)、\(l_2\)分别交于\(A\)、\(B\)，\(P(-1,2)\)为\(AB\)中点．设\(A=P+(u,v)\)，则\(B=P-(u,v)\)．\(A\in l_1\)：\(4(-1+u)+(2+v)+3=0\)，即\(4u+v+1=0\)；\(B\in l_2\)：\(3(-1-u)-5(2-v)-5=0\)，即\(-3u+5v-18=0\)．由\(v=-4u-1\)代入第二式：\(-23u-23=0\)，\(u=-1\)，\(v=3\)．故\(l\)的斜率\(k=v/u=-3\)，过\(P\)：\(y-2=-3(x+1)\)，即\(3x+y+1=0\)．（验：\(A(-2,5)\)代入\(l_1\)：\(-8+5+3=0\)，\(B(0,-1)\)代入\(l_2\)：\(0+5-5=0\)，中点为\((-1,2)=P\)，且\(A\)、\(B\)均在\(l\)上．）}
\fi
\end{ansblock}

\begin{ansblock}[2章-导-课时06-G4]
% ans:2章-导-课时06-G4
\ansitem{21}{−2}
\ifansbookdetail
\ansnote{详解}{\(x=1-2y\)即\(x+2y-1=0\)；\(2x+4y+m=0\)即\(x+2y+m/2=0\)．重合\(\iff\)两方程完全相同，故\(m/2=-1\)，\(m=-2\)．}
\fi
\end{ansblock}

% ---- 尾块（\tailfill[30mm] 定高参——钉档 §三.3：无参在平衡栏呈「内容尾随框」，贴栏底须定高参；实测无参致末页平衡 Overfull\vbox 12.99pt，定高 30mm 三零） ----
\tailfill[30mm]

\end{multicols}
\end{document}
